\documentclass{beamer}

\usepackage[utf8]{inputenc}
\input{../helper}

\title{Introdução à Computação \\ tries e hash tables}
\author{Alexandre Rademaker}
\date{}

\begin{document}

\begin{frame}
 \maketitle
\end{frame}


\begin{frame}[allowframebreaks,fragile]{trie}

  Trie (nome derivado de `retrieval') é um tipo de árvore de pesquisa
  k-ária usada para armazenar e pesquisar uma chave específica de um
  conjunto. Usando Trie, as complexidades de pesquisa pode ser
  limitada ao (comprimento da maior chave).

  Também chamadas árvores de prefixos, uma vez que prefixos comuns
  compartilham os mesmos ramos (branches).

  Podem ser usadas para armazenar um arrays associativos em que as
  chaves são normalmente cadeias de caracteres.

  \framebreak

  Imagine armazenarmos strings (chaves) em uma árvore binária de busca
  balanceada.

  Para localizar uma chave, precisariamos de um tempo proporcional a
  $M \log N$, onde M é o comprimento máximo das chaves e N é o número
  de chaves na árvore.

  Usando Trie, a chave pode ser pesquisada em tempo O(M). No entanto,
  a penalidade está nos requisitos de armazenamento do Trie.

  \framebreak

  Uma trie é uma árvore onde os nós armazenam algum valor e um array
  de ponteiros para outros nós, uma posição para cada elemento do
  alfabeto das chaves.

  \framebreak

  Uma trie com a palavra `hagrid' ficaria

  \begin{center}
    \includegraphics[width=.8\textwidth]{trie_with_hagrid.png}
  \end{center}

  \framebreak

  A medida que mais nomes forem armazenados, mais reuso de espaço é
  possível: `hagrid', `harry' e `hermione'.

  \begin{center}
    \includegraphics[width=.8\textwidth]{trie_with_multiple_names.png}
  \end{center}

\end{frame}


\begin{frame}[allowframebreaks,fragile]{trie em C}
  
  Podemos definir uma \textbf{trie} em C com

  \begin{lstlisting}[style=normalc]
    typedef struct node
    {
      bool is_word;
      struct node *children[SIZE_OF_ALPHABET];
    } t_node;
  \end{lstlisting}

  Valor \ilcode{bool} indica se este nó indica final de uma palavra
  armazenada. Mas poderíamos armazenar outros valores associados a
  chave que representa o prefixo até o dado nó.

  O array de ponteiros para os próximos nós `filhos', próximo
  caractere da chave.

  \framebreak

  A altura da \textbf{trie} é o tamanho da maior palavra armazenada.

  Independente do número de palavras, o tempo de busca de uma palavra
  será o tamanho desta palavra.

  Mas muitos espaços estarão desocupados com \ilcode{null}.
  
\end{frame}


\begin{frame}[allowframebreaks,fragile]{hash table}
  \parskip=2pt
  
  \begin{center}
    \includegraphics[width=.6\textwidth]{hash-1.png}
  \end{center}

  {\tiny \url{https://en.wikipedia.org/wiki/Hash_table}}

  \framebreak

  cada chave tem apenas um valor, mas colisões podem ocorrer

  \begin{center}
    \includegraphics[width=.6\textwidth]{hash-2.png}
  \end{center}

  {\tiny \url{https://en.wikipedia.org/wiki/Hash_table}}

  \framebreak

  Podemos pensar nas chaves como as letras A até Z, e inserir nomes de
  acordo com as iniciais.

  \begin{center}
    \includegraphics[width=.8\textwidth]{hash-az.png}
  \end{center}

  \framebreak

  E os conflitos quando tivermos que armazenar nomes com mesmas
  iniciais.

  \begin{center}
    \includegraphics[width=.8\textwidth]{hash-az-conflict.png}
  \end{center}

  \framebreak

  Cada nós terá um array de caracteres pré-allocado, de tamanho máximo
  \ilcode{LONGEST_WORD + 1} para armazenar o nome sendo guardado. O
  ponteiro \ilcode{next} aponta para outro nó, se houver conflito.

  \begin{lstlisting}[style=tinyc]
    typedef struct node
    {
      char word[LONGEST_WORD + 1];
      struct node *next;
    }
    node;
  \end{lstlisting}

  Um array de ponteiros para \ilcode{node}s com um número de `lugares'
  pré-definido, como 26, no exemplo acima (A-Z).

  \begin{lstlisting}[style=normalc]
    node *hash_table[NUMBER_OF_BUCKETS];
  \end{lstlisting}

  \framebreak

  Lembrando que a função hash mapea cada nome para uma localização
  possível da tabela hash, neste caso uma letra A-Z (ou seu código
  ASCII)

  \begin{center}
    \includegraphics[width=.55\textwidth]{hash-func.png}
  \end{center}

  \framebreak

  Os endereços na tabela hash, se associados aos códigos ASCII,
  poderiam ser traduzidos em posições diretamente.

  \begin{center}
    \includegraphics[width=.8\textwidth]{hash-az-num.png}
  \end{center}
  
  \framebreak

  Para reduzir os conflitos (e tamanho das sequências lineares de
  nós), poderíamos mudar as chaves para as \textbf{duas letras
    iniciais}
  
  \begin{center}
    \includegraphics[width=.8\textwidth]{hash-haha.png}
  \end{center}
  
  agora teríamos 676 posições na tabela, para todas as combinações das
  2 primeiras letras.

  \framebreak

  Podemos tentar aumentar o `buckets' para $26^3$ (3 iniciais)
  
  \begin{center}
    \includegraphics[width=.8\textwidth]{hash-hag.png}
  \end{center}

  reduzindo conflito mais aumentando possíveis endereços não
  ocupados. Pior caso de busca são todos os valores em uma mesma
  posição, $O(n)$, busca na lista encadeada.
  
\end{frame}




\end{document}

%%% Local Variables:
%%% mode: latex
%%% TeX-master: t
%%% End:
